% \iffalse meta-comment
%
% Copyright (C) 2025 Danila A. Kondratenko
%
% This file may be distributed and/or modified under the
% conditions of the LaTeX Project Public License, either
% version 1.3 of this license or (at your option) any later
% version. The latest version of this license is in:
%
%     http://www.latex-project.org/lppl.txt
%
% and version 1.3c or later is part of all distributions of
% LaTeX version 2008-05-04 or later.
%
% \fi
% \iffalse
%    \begin{macrocode}
%<class>\NeedsTeXFormat{LaTeX2e}
%<class>\ProvidesClass{omgtudoc-asoiu}
%<labreport>\ProvidesFile{omgtudoc-asoiu-labreport.tpt}
%<courseproj>\ProvidesFile{omgtudoc-asoiu-courseproj.tpt}
%<minimal>\ProvidesFile{omgtudoc-asoiu-minimal.tpt}
%<bibstyle>\ProvidesFile{omgtudoc-asoiu-gost-numeric.bbx}
%<*driver>
\ProvidesFile{\jobname.dtx}
%</driver>
    [2025-02-21 v1.1b
%<class>Class for ASOIU documents]
%<labreport>Laboratory report document template for omgtudoc-asoiu]
%<courseproj>ASOIU course project document template for omgtudoc-asoiu]
%<minimal>Minimal document template for omgtudoc-asoiu]
%<bibstyle>Modification of biblatex-gost gost-numeric style for omgtudoc-asoiu]
%<*driver>
]
\documentclass{ltxdoc}
\newcommand*{\rfrac}[2]{{}^{#1}\!/_{#2}}
\newcommand{\ToDo}[1]{}
\usepackage{shortvrb}
\usepackage[T2A]{fontenc}
\usepackage[utf8]{inputenc}
\usepackage[russian]{babel}
\EnableCrossrefs
\CodelineIndex
\RecordChanges
\SetupDoc{reportchangedates}
\MakeShortVerb\|
\GlossaryPrologue{\clearpage\section*{{Таблица изменений}}%
	\markboth{{Таблица изменений}}{{Таблица изменений}}}
\begin{document}
\DocInput{\jobname.dtx}
\end{document}
%</driver>
%    \end{macrocode}
% \fi
%
% ^^A \def\filedate{2025-01-16}\def\fileversion{v1.0}
% ^^A \def\filedate{2025-01-17}\def\fileversion{v1.0a}
% ^^A \def\filedate{2025-01-20}\def\fileversion{v1.0b}
% ^^A \def\filedate{2025-02-04}\def\fileversion{v1.0c}
% ^^A \def\filedate{2025-02-06}\def\fileversion{v1.0d}
% ^^A \def\filedate{2025-02-14}\def\fileversion{v1.0e}
% ^^A \def\filedate{2025-02-16}\def\fileversion{v1.1}
% ^^A \def\filedate{2025-02-18}\def\fileversion{v1.1a}
% ^^A \def\filedate{2025-02-21}\def\fileversion{v1.1b}
% \GetFileInfo{\jobname.dtx}
%
% \changes{v1.0}{2025-01-16}{Первый выпуск}
% \changes{v1.0a}{2025-01-17}{Добавлена нумерация в пределах раздела,
%    счётчики страниц, рисунков, таблиц, уравнений, приложений,
%    источников, исправлен стиль заголовков и оглавления}
% \changes{v1.0d}{2025-02-06}{Параметры документа вынесены в файлы
%    определения параметра документа}
% \changes{v1.1a}{2025-02-18}{Изменено название пакета с otuasoiu
%    на omgtudoc-asoiu}
% \changes{v1.1b}{2025-02-21}{Изменены названия файлов на более
%    локальные}
%
% \NewDocElement[macrolike=false,toplevel=false,idxtype=option,
%                idxgroup=Параметры класса,printtype=\textit{параметр}]
%               {Option}{option}
% \RenewDocElement[macrolike=false,toplevel=false,idxtype=env.,
%                idxgroup=Среда,printtype=\textit{среда}]
%               {Env}{environment}
%
% \title{\texttt{omgtudoc-asoiu}~--- класс документов кафедры АСОИУ}
%
% \author{Данила А. Кондратенко}
%
% \date{\fileversion~\filedate}
%
% \MaybeStop{^^A
%    \PrintChanges
% }
%
% \maketitle
%
% \begin{abstract}
%   Класс \texttt{omgtudoc-asoiu} предназначен для создания отчётов по
%   лабораторным, практическим работам, пояснительных записок по
%   курсовым проектам и выпускным квалификационным работам. Оформление
%   соответствует рабочей инструкции, введённой 29 января 2024 года.
% \end{abstract}
%
% \section{Введение}
%
% На кафедре АСОИУ действует стандарт оформления документов, описанный в
% соответствующей рабочей инструкции.  В ней сказано, что все документы,
% которые делают студенты, магистранты и аспиранты кафедры, должны
% подчиняться этому стандарту. Проблемой является то, что те, кто хочет
% делать документы с использованием \LaTeX, испытывают трудности в
% приведении документов в полное соответствие с этим нормативным
% документом. Данный класс позволит сократить затраты и в полной мере
% реализовать первоначальную цель \LaTeX'а~--- отделить содержимое от
% стилей.
%
% Класс |omgtudoc-asoiu| основан на стандартном классе |article|, 
% который сам по себе предоставляет практически все необходимые 
% возможности для создания документов в системе LaTeX. Проблемой его 
% является то, что форматирование этого класса не соответствует нормам
% оформления ГОСТ 7.32-2017, а именно:
%
% \begin{itemize}
% \item не соответствует шрифт;
% \item не соответствует размер листа и поля;
% \item заголовки оформляются неправильно;
% \item неправильно оформляется абзацный отступ;
% \end{itemize}
%
% И это далеко не полный список проблем с использованием стандартного
% класса |article|. Обычно эту проблему решают с применением
% дополнительных пакетов, в которых довольно сложно запутаться, в
% результате чего на оформление уходит больше времени, чем на написание
% самого отчёта.
%
% \section{Использование}
%
% Для того, чтобы использовать класс |omgtudoc-asoiu|, достаточно в преамбуле
% документа указать |\documentclass{omgtudoc-asoiu}|\relax.
%
% \subsection{Дополнительные команды}
%
% Класс |omgtudoc-asoiu| вводит следующие команды:
%
% \DescribeMacro{\appendix}
% Заголовок обязательного приложения. Принимает на вход обязательный
% аргумент \marg{name}, который указывает название приложения.
%
% \DescribeMacro{\appendixrec}
% Заголовок рекомендуемого приложения. Аналогично \cs{appendix}.
%
% \DescribeMacro{\appendixhelp}
% Заголовок справочного приложения. Аналогично \cs{appendix}.
%
% \DescribeMacro{\theappendix}
% Номер приложения.
%
% \DescribeMacro{\appendices}
% Переход в режим приложений. В этом режиме рисунки, таблицы и формулы
% нумеруются в пределах приложения.
%
% \DescribeMacro{\paragraph}
% Определение именованного пункта.
%
% \DescribeMacro{\paragraph*}
% Определение неименованного нумерованного пункта.
%
% \DescribeMacro{\subparagraph}
% Определение именованного подпункта.
%
% \DescribeMacro{\subparagraph*}
% Определение неименованного нумерованного подпункта.
%
% \DescribeMacro{\totalappendices}
% Количество приложений.
%
% \DescribeMacro{\totalfigures}
% Количество изображений.
%
% \DescribeMacro{\totaltables}
% Количество таблиц.
%
% \DescribeMacro{\totalpages}
% Общее количество страниц.
%
% \DescribeMacro{\totalsources}
% Количество источников.
%
% \subsection{Параметры}
%
% Кроме того, класс |omgtudoc-asoiu| определяет следующие типы документов:
%
% \DescribeOption{courseproj}
% Курсовой проект.
%
% \DescribeOption{labreport}
% Отчёт о лабораторной работе.
%
% \DescribeOption{minimal}
% Минимальный документ со стандартным форматированием. 
% Титульные листы не предусмотрены.
%
% Также существуют дополнительные параметры:
%
% \DescribeOption{withinsection}
% Нумерация таблиц, рисунков, формул в пределах раздела.
%
% \MaybeStop{\PrintChanges\PrintIndex}
%
% \section{Реализация}
%
% \subsection{Размеры бумаги и поля}
%
% Начнём с загрузки стандартных параметров класса |article|.
% Естественно, необходимо указать правильный размер бумаги~--- A4.
%    \begin{macrocode}
%<*class>
\LoadClass[a4paper,12pt]{article}
%    \end{macrocode}
%
% Далее укажем размер полей при помощи пакета |geometry|: левое поле~---
% 30 мм, правое~--- 15 мм, верхнее и нижнее~--- по 20 мм.
%    \begin{macrocode}
\RequirePackage[left=30mm,right=15mm,top=20mm,bottom=20mm]{geometry}
%    \end{macrocode}
%
% \subsection{Параметры}
%
% Для того, чтобы можно было выбрать подкласс документа, необходимо
% определить параметры класса.
%
% \changes{v1.1}{2025-02-16}{Добавлен тип документов minimal, он
%    установлен по умолчанию}
% \changes{v1.1b}{2025-02-21}{Изменено название файлов для типов
%    документов}
% 
%    \begin{macrocode}
\newcommand\@document@type{}
\def\CurrentOption{}
\newif\if@within@section
\@within@sectionfalse
\DeclareOption{labreport}{\renewcommand\@document@type{labreport}}
\DeclareOption{courseproj}{\renewcommand\@document@type{courseproj}}
\DeclareOption{minimal}{\renewcommand\@document@type{minimal}}
\DeclareOption{withinsection}{\@within@sectiontrue}
\ExecuteOptions{minimal}
\ProcessOptions
\input{omgtudoc-asoiu-\@document@type.tpt}
%    \end{macrocode}
%
% \subsection{Языковые настройки}
%
% Установим русский язык для создаваемого документа.
%
%    \begin{macrocode}
\RequirePackage{polyglossia}
\setdefaultlanguage[forceheadingpunctuation=false]{russian}
%    \end{macrocode}
%
% Далее нужно установить формат подписей к рисункам и таблицам. Подпись
% состоит из обозначения объекта (таблица или рисунок) и длинного тире
% между ними.
%
%    \begin{macrocode}
\addto\captionsrussian{%
    \renewcommand\figurename{Рисунок}%
    \renewcommand\tablename{Таблица}%
}
\RequirePackage{caption}
\DeclareCaptionLabelSeparator{emdash}{~---~}
\captionsetup[figure]{labelsep=emdash,justification=centerlast}
\captionsetup[table]{labelsep=emdash,justification=justified,singlelinecheck=false}
%    \end{macrocode}
%
% \subsection{Общее форматирование текста}
%
% Укажем шрифт Tempora, который похож на Times New Roman, для обычного
% текста, и шрифт из пакета \texttt{newtxmath}.
%
% \changes{v1.0e}{2025-02-14}{Исправлены проблемы с лигатурами \TeX\ и определением
%     кириллицы в шрифтах: шрифты заменены на Tempora и newtxmath}
%
%    \begin{macrocode}
\renewcommand*{\rmdefault}{minntx}
\RequirePackage[]{newtxmath}
\RequirePackage{tempora}
\RequirePackage[no-math]{fontspec}
\setmainfont{tempora}
%    \end{macrocode}
%
% Установим шрифт в 14 пт и прочие соответствующие шрифты.
%
% \changes{v1.0b}{2025-01-20}{Исправлены настройки шрифтов}
%
%    \begin{macrocode}
\renewcommand\normalsize{\@setfontsize\normalsize{14}{16}%
    \abovedisplayskip 14pt plus 3pt minus 7pt%
    \abovedisplayshortskip 0pt plus 4pt%
    \belowdisplayshortskip 7pt plus 4pt minus 3pt%
    \belowdisplayskip \abovedisplayskip}
\renewcommand\small{\@setfontsize\small\@xiipt{14}%
    \abovedisplayskip 12pt plus 3pt minus 7pt%
    \abovedisplayshortskip 0pt plus 3pt%
    \belowdisplayshortskip 6pt plus 3pt minus 3pt%
    \belowdisplayskip \abovedisplayskip}
\renewcommand\footnotesize{\@setfontsize\footnotesize\@xpt\@xiipt}
\renewcommand\scriptsize{\@setfontsize\scriptsize\@viiipt{9.5}}
\renewcommand\tiny{\@setfontsize\tiny\@vipt\@viipt}
\renewcommand\large{\@setfontsize\large\@xviipt{22}}
\renewcommand\Large{\@setfontsize\Large\@xxpt{25}}
\renewcommand\LARGE{\@setfontsize\LARGE\@xxvpt{30}}
\renewcommand\huge{\@setfontsize\huge{30}{35}}
\renewcommand\Huge{\@setfontsize\Huge{36}{40}}
\normalsize
%    \end{macrocode}
%
% \begin{macro}{\linespread}
% Установим межстрочный интервал в 130%.
%    \begin{macrocode}
\linespread{1.3}
%    \end{macrocode}
% \end{macro}
%
% \begin{macro}{\parskip}
% Дополнительных промежутков между абзацами быть не должно.
%    \begin{macrocode}
\setlength\parskip{0pt}
%    \end{macrocode}
% \end{macro}
%
% \begin{macro}{\parindent}
% Установим абзацный отступ в 12,5 мм.
%    \begin{macrocode}
\setlength\parindent{12.5mm}
%    \end{macrocode}
% \end{macro}
%
% \begin{macro}{\ps@plain}
% Зададим стиль страниц по умолчанию (это необходимо, чтобы на всех страницах
% номер писался одним и тем же шрифтом):
%    \begin{macrocode}
\def\ps@plain{\let\@mkboth\@gobbletwo
     \let\@oddhead\@empty\def\@oddfoot{\hfil\rmfamily\thepage
     \hfil}\let\@evenhead\@empty\let\@evenfoot\@oddfoot}
\pagestyle{plain}
%    \end{macrocode}
% \end{macro}
%
% \subsection{Заголовки}
%
% Изменим стиль заголовков первых двух уровней. Разделы должны
% начинаться с новой страницы, сами заголовки от остальных абзацев
% отличаются не сильно, каждый заголовок раздела начинается с абзацного
% отступа, за которым идёт номер и название раздела.
%    \begin{macrocode}
\renewcommand\section{\clearpage\@startsection{section}{1}{\parindent}%
        {0pt}{1sp}{\normalfont\normalsize\bfseries}}
\renewcommand\subsection{\@startsection{subsection}{2}{\parindent}%
    {0pt}{1sp}{\normalfont\normalsize\bfseries}}
%    \end{macrocode}
%
% \subsection{Переопределение модели структуры документа}
%
% Кафедральная инструкция определяет, что существует четыре уровня
% дробления разделов:
%
% \begin{itemize}
% \item разделы;
% \item подразделы;
% \item пункты;
% \item подпункты.
% \end{itemize}
%
%    \begin{macrocode}
\setcounter{secnumdepth}{4}
%    \end{macrocode}
%
% Определим команды \cs{paragraph} и \cs{paragraph*},
% \cs{subparagraph} и \cs{subparagraph*}.
%
% \begin{macro}{\paragraph}\oarg{tocname}\marg{name}. Записывает пункт в
% содержание под названием \marg{name} или \oarg{tocname}, если имеется,
% и выдаёт пункт с названием, выделенным полужирным шрифтом.
%    \begin{macrocode}
\renewcommand\paragraph{\@startparagraph{paragraph}{3}}
%    \end{macrocode}
% \end{macro}
%
% \begin{macro}{\paragraph*} Выдаёт неназванный нумерованный пункт.
% \end{macro}
%
% \begin{macro}{\subparagraph}\oarg{tocname}\marg{name}. Выдаёт подпункт
% с заголовком, выделенным жирным шрифтом.
%    \begin{macrocode}
\renewcommand\subparagraph{\@startparagraph{subparagraph}{4}}
%    \end{macrocode}
% \end{macro}
%
% \begin{macro}{\subparagraph*} Выдаёт неназванный нумерованный
% подпункт.
% \end{macro}
%
% Поскольку эти команды имеют много общего, определим общую команду
% \cs{@startparagraph}:
%
% \begin{macro}{\@startparagraph}\marg{name}\marg{level}*\oarg{tocname}\marg{name}.
%    \begin{macrocode}
\def\@startparagraph#1#2{
    \if #2>\c@secnumdepth%
    \else%
    \refstepcounter{#1}%
    \fi%
    \@ifstar{{}\unnamed@paragraph{#1}{#2}}{\@dblarg{\named@paragraph{#1}{#2}}}%
}
%    \end{macrocode}
% \end{macro}
%
% \begin{macro}{\named@paragraph}\marg{name}\marg{level}\oarg{tocname}\marg{name}
%    \begin{macrocode}
\def\named@paragraph#1#2[#3]#4{%
    \addcontentsline{toc}{#1}{%
        \protect\numberline{%
            \csname the#1\endcsname\space}%
            #3%
    }%\par
    \def\otu@partitle{{#4}}%
    \paragraph@heading{#1}%
}
%    \end{macrocode}
% \end{macro}
% \begin{macro}{\unnamed@paragraph}\marg{name}\marg{level}
%    \begin{macrocode}
\def\unnamed@paragraph#1#2{
    \let\otu@partitle\@empty%
    \paragraph@heading{#1}%
}
%    \end{macrocode}
% \end{macro}
%
% Теперь нужно будет сделать так, чтобы название пункта/подпункта соединялось
% со следующим абзацем. Код был заимствован из стандартной поставки \LaTeX
% и изначально отвечал за печать заголовков со значением отступа после
% заголовка меньше 0.
%
% \begin{macro}{\paragraph@heading}\marg{name}
% \changes{v1.0d}{2025-02-06}{Исправлено размещение пробела между номером пункта
%     и остальным текстом}
%    \begin{macrocode}
\def\paragraph@heading#1{%
    \@nobreakfalse
    \global\@noskipsectrue
    \everypar{%
      \if@noskipsec
        \global\@noskipsecfalse
       {\setbox\z@\lastbox}%
        \clubpenalty\@M
        \begingroup%
            \hskip\parindent%
            \normalfont\normalsize\bfseries
            \csname the#1\endcsname\,%
            \ifx\otu@partitle\empty\relax\else\otu@partitle\relax\fi\,%
        \endgroup
        \unskip
      \else
        \clubpenalty \@clubpenalty
        \everypar{}%
      \fi}%
    \ignorespaces%
}
%    \end{macrocode}
% \end{macro}
%
% Далее переопределим связи между счётчиками уровней.
%
%    \begin{macrocode}
\counterwithout{paragraph}{subsubsection}
\counterwithout{subsubsection}{subsection}
\counterwithout{subsubsection}{section}
\counterwithin{paragraph}{subsection}
%    \end{macrocode}
%
% Далее определим формат нумерации для всех уровней. Разделы,
% подразделы, пункты и подпункты нумеруются арабскими цифрами, все цифры
% объединяются через точки.
%
%    \begin{macrocode}
\renewcommand\theparagraph{\thesubsection.\@arabic\c@paragraph}
\renewcommand\thesubparagraph{\theparagraph.\@arabic\c@subparagraph}
%    \end{macrocode}
%
% В стандартном классе \LaTeX{} между пунктами и подпунктами имеется ещё
% одна степень дробления~--- подподразделы. Её необходимо исключить:
%
%    \begin{macrocode}
\let\subsubsection\@@@undefined@@@
\let\l@subsubsection\@@@undefined@@@
\let\c@subsubsection\@@@undefined@@@
%    \end{macrocode}
%
% \subsection{Содержание}
%
% Теперь приступим к форматированию содержания. В \LaTeX за это отвечают
% команды вида \cs{l@}\meta{level}. Необходимо, чтобы в содержании после
% заголовков каждого элемента стояло отточие и чтобы соблюдались отступы
% подразделов и пунктов в 2 и 4 знака соответственно.
%
% \begin{macro}{\l@section,\l@subsection,\l@paragraph}
%    \begin{macrocode}
\renewcommand*\l@section[2]{\@dottedtocline{1}{0em}{1em}{\textbf{#1}}{#2}}
\renewcommand*\l@subsection{\@dottedtocline{2}{2ex}{2em}}
\renewcommand*\l@paragraph{\@dottedtocline{3}{4ex}{3em}}
%</class>
%    \end{macrocode}
% \end{macro}
%
% \subsection{Библиография}
%
% Для библиографии будем использовать пакет |biblatex| и
% |biblatex-gost|.
%
% Для начала исправим один изъян в генерации библиографии.
% Библиографические записи должны начинаться с абзацного отступа и
% нумероваться. Для того, чтобы это реализовать, необходимо исправить
% стилевой файл библиографии:
%
%    \begin{macrocode}
%<*bibstyle>
\RequireBibliographyStyle{gost-standard}
\ExecuteBibliographyOptions{labelnumber,bibisbn,biburl,bibdoi,bibeprint}

\providebool{bbx:subentry}

\DeclareBibliographyOption{subentry}[true]{%
  \setbool{bbx:subentry}{#1}}

\DeclareFieldFormat{bibentrysetcount}{\mkbibparens{\mknumalph{#1}}}
\DeclareFieldFormat{labelnumberwidth}{#1\addperiod}
\DeclareFieldFormat{shorthandwidth}{\mkbibbrackets{#1}}

\defbibenvironment{bibliography}
  {\list
     {\printtext[labelnumberwidth]{%
    \printfield{labelprefix}%
    \printfield{labelnumber}}}
     {\setlength{\labelwidth}{0pt}%
      \setlength{\leftmargin}{0pt}%
      \setlength{\labelsep}{0pt}%
      \setlength{\itemindent}{\parindent}%
      \setlength{\itemsep}{0pt}%
      \setlength{\parsep}{0pt}}%
%    \end{macrocode}
% В этой строке устанавливается формат номера в библиографии. Такая же
% модификация сделана для короткой записи библиографии (|shorthand|).
%    \begin{macrocode}
    \renewcommand*{\makelabel}[1]{##1\>}}
  {\endlist}
  {\item}

\defbibenvironment{shorthand}
  {\list
     {\printfield[shorthandwidth]{shorthand}}
     {\setlength{\labelwidth}{\shorthandwidth}%
      \setlength{\leftmargin}{\labelwidth}%
      \setlength{\labelsep}{\biblabelsep}%
      \addtolength{\leftmargin}{\labelsep}%
      \setlength{\itemsep}{\bibitemsep}%
      \setlength{\parsep}{\bibparsep}%
      \renewcommand*{\makelabel}[1]{\hspace{12.5mm}\hss##1\>}}}
  {\endlist}
  {\item}

\DeclareBibliographyDriver{set}{%
  \entryset
    {\ifbool{bbx:subentry}
       {\printfield[bibentrysetcount]{entrysetcount}%
    \printunit*{\addnbspace}}
       {}}
    {}%
  \newunit\newblock
  \usebibmacro{setpageref}%
  \finentry}
%</bibstyle>
%    \end{macrocode}
%
% \begin{macro}{biblatex}
% Теперь можно подключать пакет |biblatex|:
% \changes{v1.1b}{2025-02-21}{Изменено название стиля библиографии}
%    \begin{macrocode}
%<*class>
\RequirePackage[sorting=none,bibstyle=omgtudoc-asoiu-gost-numeric,%
citestyle=gost-numeric]{biblatex}
%    \end{macrocode}
% \end{macro}
%
% \subsection{Заголовки структурных частей документа}
% 
% Стандарт предполагает следующие структурные части:
% 
% \begin{itemize}
% \item содержание;
% \item введение;
% \item перечень сокращений и обозначений;
% \item термины и определения (может быть объединена с предыдущим
% пунктом под названием <<определения, сокращения и обозначения>>;
% \item заключение;
% \item список использованных источников;
% \item реферат;
% \item список исполнителей.
% \end{itemize}
%
% Заголовки этих структурных частей размещаются посередине. Реализуем
% команды для их помещения.
%
% Для начала реализуем общую команду для подобного рода заголовков
% \cs{centered@heading}, которая будет помещать ненумерованный
% отцентрованный заголовок в текст и в содержание.
%
% \begin{macro}{\centered@heading}
% \changes{v1.0d}{2025-02-06}{Исправлено выравнивание заголовка
%    отцентрованного ненумерованного с помещением в содержание}
% Заголовок отцентрованный ненумерованный, помещается на уровень раздела
% в содержание.
%    \begin{macrocode}
\def\centered@heading#1{%
    \clearpage%
    \par%
    \addcontentsline{toc}{section}{#1}%
    \begingroup%
    \centering\normalfont\normalsize\bfseries\MakeUppercase%
    {\hfill #1 \hfill}
    \endgroup \par
}
%    \end{macrocode}
% \end{macro}
%
% \begin{macro}{\centered@heading@nontoc}
% \changes{v1.0d}{2025-02-06}{Исправлено выравнивание заголовка
%    отцентрованного ненумерованного без помещения в содержание}
% Заголовок отцентрованный ненумерованный, в содержание не помещается.
%    \begin{macrocode}
\def\centered@heading@nontoc#1{%
    \begingroup%
    \centering\normalfont\normalsize\bfseries\MakeUppercase%
    {\hfill #1 \hfill}
    \endgroup \par
}
%    \end{macrocode}
% \end{macro}
% 
% Далее~--- определим уже сами макросы.
%
% \begin{macro}{\preface} Заголовок введения.
%    \begin{macrocode}
\newcommand{\preface}{\centered@heading{Введение}}
%    \end{macrocode}
% \end{macro}
%
% \begin{macro}{\conclusion} Заголовок заключения.
%    \begin{macrocode}
\newcommand{\conclusion}{\centered@heading{Заключение}}
%    \end{macrocode}
% \end{macro}
%
% \begin{macro}{\abbreviations} Заголовок перечня сокращений и
% обозначений.
%    \begin{macrocode}
\newcommand{\abbreviations}{\centered@heading{Перечень сокращений и
обозначений}}
%    \end{macrocode}
% \end{macro}
%
% \begin{macro}{\terms} Заголовок списка терминов и определений.
%    \begin{macrocode}
\newcommand{\terms}{\centered@heading{Термины и определения}}
%    \end{macrocode}
% \end{macro}
%
% \begin{macro}{\termsandabbrev} Заголовок списка терминов и
% сокращений.
%    \begin{macrocode}
\newcommand{\termsandabbrev}{\centered@heading{Определения,
сокращения и обозначения}}
%    \end{macrocode}
% \end{macro}
%
% Определим также заголовок содержания. Параметры его находятся в
% команде \cs{tableofcontents}, которую необходимо переопределить:
%
% \changes{v1.0d}{2025-02-06}{Убраны избыточные команды из определения оглавления}
%
%    \begin{macrocode}
\renewcommand\tableofcontents{%
    \centered@heading@nontoc{\contentsname}%
    \@starttoc{toc}%
    }
%    \end{macrocode}
%
% Определим также заголовок для списка использованных источников:
% \begin{macro}{otuasoiubibliography}
% \changes{v1.0d}{2025-02-06}{Убраны избыточные команды из определения
%     заголовка библиографии}
% Стиль заголовка для списка использованных источников в соответствии с
% ГОСТ 7.32-2017 и рабочей инструкцией.
%    \begin{macrocode}
\defbibheading{otuasoiubibliography}[\bibname]{%
    \centered@heading{#1}}
\DeclarePrintbibliographyDefaults{heading=otuasoiubibliography}
%    \end{macrocode}
% \end{macro}
%
% \subsection{Формулы}
%
% Согласно инструкции, формулы обязательно нумеруются, сама формула
% помещается по центру, а номер формулы помещается справа. Для того,
% чтобы добиться такого вида, необходимо использовать пакет |amsmath|.
%
%    \begin{macrocode}
\RequirePackage{amsmath}
%    \end{macrocode}
%
% Ссылки на формулы указываются в скобках: для этого используется команда
% \cs{eqref}.
%
% \subsection{Нумерация}
%
% Нумерация рисунков, таблиц и формул может осуществляться как по всему
% документу, так и по разделам. Для установки нумерации по разделам
% применяется параметр |withinsection|.
%
%    \begin{macrocode}
\if@within@section
\numberwithin{equation}{section}
\counterwithin{table}{section}
\counterwithin{figure}{section}
\fi
%    \end{macrocode}
%
% \subsection{Счётчики}
%
% Несмотря на введение |withinsection|, необходимо считать количество
% таблиц и рисунков. Тем не менее, при установке счётчиков в 0
% информация о количестве может потеряться.
%
%    \begin{macrocode}
\RequirePackage{totcount}
\RequirePackage{assoccnt}
\RequirePackage{xspace}
%    \end{macrocode}
%
% \begin{macro}{\totalfigures} Количество рисунков.
%    \begin{macrocode}
\newtotcounter{@figures}
\regtotcounter{figure}
\DeclareAssociatedCounters{figure}{@figures}
\newcommand\totalfigures{\total{@figures}\xspace}
%    \end{macrocode}
% \end{macro}
%
% \begin{macro}{\totaltables} Количество таблиц.
%    \begin{macrocode}
\newtotcounter{@tables}
\regtotcounter{table}
\DeclareAssociatedCounters{table}{@tables}
\newcommand\totaltables{\total{@tables}\xspace}
%    \end{macrocode}
% \end{macro}
%
% \begin{macro}{\totalappendices} Количество приложений.
%    \begin{macrocode}
\newtotcounter{@appendices}
\setcounter{@appendices}{0}
\DeclareAssociatedCounters{appendix}{@appendices}
\newcommand\totalappendices{\total{@appendices}\xspace}
%    \end{macrocode}
% \end{macro}
%
% \begin{macro}{\totalequations} Количество уравнений.
%    \begin{macrocode}
\newtotcounter{@equations}
\DeclareAssociatedCounters{equation}{@equations}
\newcommand\totalequations{\total{@equations}\xspace}
%    \end{macrocode}
% \end{macro}
%
% \begin{macro}{\totalpages} Количество страниц.
%    \begin{macrocode}
\regtotcounter{page}
\newcommand\totalpages{\total{page}\xspace}
%    \end{macrocode}
% \end{macro}
%
% \begin{macro}{\totalsources} Количество источников.
%    \begin{macrocode}
\newtotcounter{@sourceref}
\AtEveryBibitem{\stepcounter{@sourceref}}
\newcommand\totalsources{\total{@sourceref}\xspace}
%    \end{macrocode}
% \end{macro}
%
% \subsection{Приложения}
%
% Теперь реализуем приложения. Для отделения основного текста от
% приложений будем использовать команду \cs{appendices}, которая
% переопределяет некоторые счётчики.
%
% Для нумерации приложений будут использоваться заглавные буквы русского
% алфавита за исключением Ё, З, Й, О, Ч, Ъ, Ы, Ь. З, О и Ч похожи на
% цифры 3, 0 и 4. Доступ к этим буквам можно будет получить через
% команду \cs{GostAsbuk}.
%
%    \begin{macrocode}
\def\gostasbuk@Alph#1{\ifcase#1\or
   А\or Б\or В\or Г\or Д\or Е\or Ж\or
   И\or К\or Л\or М\or Н\or П\or Р\or
   С\or Т\or У\or Ф\or Х\or Ц\or Ш\or 
   Щ\or Э\or Ю\or Я%
   \else\ClassError{omgtudoc-asoiu}{Invalid value of counter}\fi
}
\def\GostAsbuk#1{\expandafter\gostasbuk@Alph\csname c@#1\endcsname}
%    \end{macrocode}
%
% Теперь можно реализовать команду \cs{appendices}, которая переведёт
% дальнейший документ в режим приложений. В этом режиме в оглавление
% будут попадать только заголовки самих приложений, но не их
% подразделов.
%
% \begin{macro}{\appendices}
%    \begin{macrocode}
\newtotcounter{appendix}
\newcommand\appendices{%
    \counterwithin{section}{appendix}%
    \addtocontents{toc}{\setcounter{tocdepth}{0}}%
    \renewcommand\thesection{\theappendix.\arabic{section}}%
    \renewcommand\section{\@startsection{section}{1}{\parindent}%
        {0pt}{1sp}{\normalfont\normalsize\bfseries}}%
    \renewcommand\theappendix{\GostAsbuk{appendix}}%
%    \end{macrocode}
%
% Далее нужно будет переопределить счётчики.
%
%    \begin{macrocode}
    \counterwithout{equation}{section}%
    \counterwithout{figure}{section}%
    \counterwithout{table}{section}%
    \numberwithin{equation}{appendix}%
    \counterwithin{figure}{appendix}%
    \counterwithin{table}{appendix}
}
%    \end{macrocode}
% \end{macro}
%
% Теперь приложения начинаются с команд вида
% \cs{appendix}\meta{type}\oarg{tocname}\marg{name}. Существует три типа
% приложений:
%
% \begin{itemize}
% \item обязательное (команда \cs{appendix});
% \item рекомендуемое (команда \cs{appendixrec});
% \item справочное (команда \cs{appendixhelp})
% \end{itemize}
%
% Эти три команды имеют общую структру, поэтому они будут вызывать более
% общую команду \cs{@startappendix}:
%
% \begin{macro}{\@startappendix}\marg{type}\marg{name}
%
% Печатает заголовок приложения с типом \meta{type} и названием
% \meta{name}.
%
%    \begin{macrocode}
\def\@startappendix#1#2{
    \clearpage%
    \SuspendCounters{@figures}%
    \SuspendCounters{@tables}%
    \SuspendCounters{@equations}%
    \refstepcounter{appendix}%
    \ResumeSuspendedCounters%
    \addcontentsline{toc}{appendix}{Приложение \theappendix\space (#1) #2}%
    \begingroup
    \centering\normalfont\normalsize\bfseries%
    ПРИЛОЖЕНИЕ \theappendix{}\\
    \centering ({#1})\\
    \centering {\hfill #2 \hfill}
    \endgroup
    \par
}
%    \end{macrocode}
% \end{macro}
% \begin{macro}{\appendix,\appendixhelp,\appendixrec}
%    \begin{macrocode}
\renewcommand{\appendix}{\@startappendix{обязательное}}
\newcommand{\appendixhelp}{\@startappendix{справочное}}
\newcommand{\appendixrec}{\@startappendix{рекомендуемое}}
%    \end{macrocode}
% \end{macro}
%
% Определим также стиль пункта оглавления для приложений.
%
%    \begin{macrocode}
\newcommand*\l@appendix[2]{\@dottedtocline{0}{0em}{3em}{\textbf{#1}}{#2}}
%    \end{macrocode}
%
% \subsection{Титульные листы} \label{titlepage}
%
% Первой страницей любого отчёта является титульный лист, который не
% нумеруется и содержит все основные данные о работе.
%
% \begin{macro}{\maketitle} Создаёт титульный лист.
%
% \changes{v1.1}{2025-02-16}{Титульный лист не создаётся, если не определена
%    команда \cs{@titlepage@template}}
%
%    \begin{macrocode}
\RequirePackage{etoolbox}
\renewcommand\maketitle{%
    \ifcsdef{@titlepage@template}{\@titlepage@template%
    \thispagestyle{empty}%
    \newpage}{}
}
%</class>
%    \end{macrocode}
% \end{macro}
%
% \begin{macro}{\@titlepage@template}
% Команда \cs{@titlepage@template} определяется в одном из файлов с
% расширением |tpt|, который соответствует переданному параметру.
%
%    \begin{macrocode}
%<*labreport>
\long\def\@titlepage@template{%
    \begin{center}%
        Министерство науки и высшего образования Российской Федерации

        Федеральное государственное автономное образовательное
        учреждение высшего образования

        <<Омский государственный технический университет>>

        Кафедра <<Автоматизированные системы обработки информации и
        управления>>
        ~\par
        ~\par
        ~\par
        \textbf{ОТЧЁТ}\par
        о лабораторной работе\@laboratory@index\@laboratory@theme
        по дисциплине
        <<\@laboratory@subject>>\par
        студента \@student@surname@g\ \@student@name@g\space
        \@student@fathername@g\ группы \@student@group\par
        ~\par
        ~\par
        \textbf{Пояснительная записка}\par
        Шифр проекта ЛР-02068999-43-\@student@id-ПЗ\par
        \@student@direction@a\ \@student@direction@id\space
        \@student@direction@name\par
        \vfill
        \@professor@value \hfill \@professor@initials\par
        ~\par
        Студент \hfill \@student@initials\par
        \vfill
        Омск \@doc@year
    \end{center}
}
%</labreport>
%<*courseproj>
\long\def\@titlepage@template{%
    \begin{center}
        Министерство науки и высшего образования Российской Федерации

        Федеральное государственное автономное образовательное
        учреждение высшего образования

        <<Омский государственный технический университет>>

        Кафедра <<Автоматизированные системы обработки информации и
        управления>>
        ~\par
        ~\par
        ~\par
        {\bfseries КУРСОВОЙ ПРОЕКТ}\par
        на тему <<\@courseproj@theme>>\par
        по дисциплине <<\@courseproj@subject>>\par
        студента \@student@surname@g\ \@student@name@g\space
        \@student@fathername@g\ группы \@student@group\par
        ~\par
        ~\par
        {\bfseries Пояснительная записка}\par
        Шифр проекта КП-02068999-43-\@student@id-ПЗ\par
        \@student@direction@a\space \@student@direction@id\space
            \@student@direction@name\par
        \vfill
        \@professor@value\hfill \@professor@initials\par
        ~\par
        Студент \hfill \@student@initials\par
        ~\par
        {\raggedleft К защите:\hfill}
        \vskip\baselineskip
        \begin{tabular}{|c|c|c|}
            \hline
            \begin{minipage}[t]{5.5cm}
            \begin{center}
            \setlength{\parindent}{0em}
            Выполнение и подготовка к защите, баллы
            \end{center}
            \end{minipage} & 
            \begin{minipage}[t]{5cm}
            \begin{center}
            Защита КП, баллы
            \end{center} 
            \end{minipage} & 
            \begin{minipage}[t]{5cm}
            \begin{center}
            Итоговый рейтинг, баллы
            \end{center}
            \end{minipage} \\
            \hline
            ~&~&~ \\ \hline
        \end{tabular}\par
        \vskip\baselineskip
    \end{center}
    \begin{flushleft}
        {Проект защищён с оценкой:}\par
    \end{flushleft}
    \vfill
    \begin{center}
        Омск \@doc@year
        \setlength{\parindent}{12.5mm}
    \end{center}
}
%</courseproj>
%    \end{macrocode}
%
% Для типа документа |minimal| титульный лист не предусмотрен.
%    \begin{macrocode}
%<*minimal>
\ClassWarning{omgtudoc-asoiu}{Minimal document type is selected: %
   no metadata setting commands are available, title page %
   won't be created.}
%</minimal>
%    \end{macrocode}
% \end{macro}
%
% \subsection{Параметры документа}
%
% Класс |otuasoiu| определяет следующие общие параметры:
%
% \begin{macro}{\studentname}\marg{surnamen}\marg{namen}
% \marg{fathernamen}\marg{surnameg}\marg{nameg}\marg{fathernameg}.
% 
% Определяет имя студента в именительном и родительном падежах.
%
% \changes{v1.0c}{2025-02-04}{Исправлено получение инициалов студента}
% \changes{v1.0d}{2025-02-06}{Все параметры студента (ФИО) перемещены
%    ближе к определению \cs{studentname}}
% \ToDo{вынести получение инициалов в отдельную функцию}
% \ToDo{реализовать склонение для имён - склонение вынести в отдельный пакет}
%
%    \begin{macrocode}
%<*labreport|courseproj>
\gdef\@student@surname@n{\ClassError{omgtudoc-asoiu}{Не указана фамилия
студента}}
\gdef\@student@surname@g{\ClassError{omgtudoc-asoiu}{Не указана фамилия
студента}}
\gdef\@student@name@n{\ClassError{omgtudoc-asoiu}{Не указано имя
студента}}
\gdef\@student@name@g{\ClassError{omgtudoc-asoiu}{Не указано имя
студента}}
\gdef\@student@fathername@n{\ClassError{omgtudoc-asoiu}{Не указано отчество
студента}}
\gdef\@student@fathername@g{\ClassError{omgtudoc-asoiu}{Не указано отчество
студента}}
\ExplSyntaxOn
\str_new:N \l__otuasoiu_nameinitial_str
\str_new:N \l__otuasoiu_fathernameinitial_str
\tl_new:N \l__otuasoiu_shortname_tl
\newcommand{\studentname}[6]{%
    \xdef\@student@surname@n{#1}%
    \xdef\@student@name@n{#2}%
    \xdef\@student@fathername@n{#3}%
    \xdef\@student@surname@g{#4}%
    \xdef\@student@name@g{#5}%
    \xdef\@student@fathername@g{#6}%
    \str_set:Ne \l__otuasoiu_nameinitial_str 
        { \str_item:Nn { \@student@name@n } {1} }
    \str_set:Ne \l__otuasoiu_fathernameinitial_str 
        { \str_item:Nn { \@student@fathername@n } {1} }
    \tl_set:Nn \l__otuasoiu_shortname_tl { 
        \text_uppercase:n { \str_use:N \l__otuasoiu_nameinitial_str } .~ 
        \text_uppercase:n { \str_use:N \l__otuasoiu_fathernameinitial_str } .~
        \@student@surname@n }
    \xdef\@student@initials{\tl_use:N \l__otuasoiu_shortname_tl}
}
\ExplSyntaxOff
%</labreport|courseproj>
%    \end{macrocode}
% \end{macro}
%
% \begin{macro}{\studentgroup}\marg{group}.
%
% Определяет группу студента.
%
%    \begin{macrocode}
%<*labreport|courseproj>
\gdef\@student@group{\ClassError{omgtudoc-asoiu}{Не указана группа студента}}
\newcommand{\studentgroup}[1]{\gdef\@student@group{#1}}
%</labreport|courseproj>
%    \end{macrocode}
% \end{macro}
%
% \begin{macro}{\studentid}\marg{id}.
%
% Определяет номер студента.
%
%    \begin{macrocode}
%<*labreport|courseproj>
\gdef\@student@id{\ClassError{omgtudoc-asoiu}{Не указан номер зачётной
книжки/студенческого билета}}
\newcommand{\studentid}[1]{\gdef\@student@id{#1}}
%</labreport|courseproj>
%    \end{macrocode}
% \end{macro}
%
% \begin{macro}{\speciality}\marg{spec}\marg{id}\marg{name}.
%
% Определяет специальность студента. В поле \marg{spec} вписывается обычно или
% <<Направление>>, или <<Специальность>>.
%
%    \begin{macrocode}
%<*labreport|courseproj>
\gdef\@student@direction@a{\ClassError{omgtudoc-asoiu}{Не указано
направление студента}}
\gdef\@student@direction@id{\ClassError{omgtudoc-asoiu}{Не указано
направление студента}}
\gdef\@student@direction@name{\ClassError{omgtudoc-asoiu}{Не указано
направление студента}}
\newcommand{\speciality}[3]{%
    \gdef\@student@direction@a{#1}%
    \gdef\@student@direction@id{#2}%
    \gdef\@student@direction@name{#3}%
}
%    \end{macrocode}
% \end{macro}
%
% \changes{v1.1}{2025-02-16}{Добавление макросов для конкретных значений направления:
%    \cs{directionIVT}, \cs{directionPIN}, \cs{specialityPE}}
%
% \begin{macro}{\directionIVT}
% Задаёт направление 09.03.01 <<Информатика и вычислительная техника>>
%    \begin{macrocode}
\newcommand{\directionIVT}{\speciality{Направление}%
	{09.03.01}{Информатика и вычислительная техника}}
%    \end{macrocode}
% \end{macro}
% \begin{macro}{\directionPIN}
% Задаёт направление 09.03.04 <<Программная инженерия>>
%    \begin{macrocode}
\newcommand{\directionPIN}{\speciality{Направление}%
	{09.03.04}{Программная инженерия}}
%    \end{macrocode}
% \end{macro}
% \begin{macro}{\specialityPE}
% \changes{v1.1b}{2025-02-21}{Исправлено определение макроса \cs{specialityPE}}
% Задаёт специальность 09.05.01 <<Применение и эксплуатация автоматизированных
% систем специального назначения>>.
%    \begin{macrocode}
\newcommand{\specialityPE}{\speciality{Специальность}%
	{09.05.01}{Применение и эксплуатация автоматизированных систем %
	специального назначения}}
%</labreport|courseproj>
%    \end{macrocode}
% \end{macro}
%
% \begin{macro}{\laboratoryid}\marg{id}.
%
% Определяет номер лабораторной работы. Только для типа документа |labreport|.
%
%    \begin{macrocode}
%<*labreport>
\gdef\@laboratory@index{\ClassWarning{omgtudoc-asoiu}{Не указан номер
лабораторной работы}} 
\ExplSyntaxOn
\newcommand{\laboratoryid}[1]{\gdef\@laboratory@index{~№#1~}}
\ExplSyntaxOff
%</labreport>
%    \end{macrocode}
% \end{macro}
%
% \begin{macro}{\laboratorytheme}{\marg{name}}.
%
% Определяет тему лабораторной работы. Только для типа документа |labreport|.
%
%    \begin{macrocode}
%<*labreport>
\gdef\@laboratory@theme{\ClassWarning{omgtudoc-asoiu}{Не указана тема
лабораторной работы}\par}
\newcommand{\laboratorytheme}[1]{\gdef\@laboratory@theme{\par<<#1>>\par}}
%</labreport>
%    \end{macrocode}
% \end{macro}
%
% \begin{macro}{\laboratorysubject}\marg{name}.
%
% Определяет дисциплину, по которой выполнена лабораторная работа. Только для
% типа документа |labreport|.
%
% \changes{v1.0e}{2025-02-14}{Исправлена ошибка: непонятный 0 в
%     определении}
%    \begin{macrocode}
%<*labreport>
\gdef\@laborartory@subject{\ClassError{omgtudoc-asoiu}{Не указана
дисциплина, по которой выполнена лабораторная работа}}
\newcommand{\laboratorysubject}[1]{\gdef\@laboratory@subject{#1}}
%</labreport>
%    \end{macrocode}
% \end{macro}
%
% \begin{macro}{\courseprojsubject}\marg{name}.
% 
% Дисциплина, по которой сдаётся курсовой проект. Только для типа документа
% |courseproj|.
%
%    \begin{macrocode}
%<*courseproj>
\gdef\@courseproj@subject{\ClassError{omgtudoc-asoiu}{Не указана дисциплина,
по которой выполнен курсовой проект}}
\newcommand{\courseprojsubject}[1]{\gdef\@courseproj@subject{#1}}
%</courseproj>
%    \end{macrocode}
% \end{macro}
%
% \begin{macro}{\courseprojtheme}\marg{name}.
% 
% Тема курсового проекта. Только для типа документа |courseproj|.
%
%    \begin{macrocode}
%<*courseproj>
\gdef\@courseproj@theme{\ClassError{omgtudoc-asoiu}{Не указана тема курсового
проекта}}
\newcommand{\courseprojtheme}[1]{\gdef\@courseproj@theme{#1}}
%</courseproj>
%    \end{macrocode}
% \end{macro}
%
% \begin{macro}{\professorvalue}\marg{value}.
%
% Должность и степень преподавателя.
%
%    \begin{macrocode}
%<*labreport|courseproj>
\gdef\@professor@value{\ClassError{omgtudoc-asoiu}{Не указан
преподаватель}}
\gdef\@professor@initials{\ClassError{omgtudoc-asoiu}{Не указан
преподаватель}}
\newcommand{\professorvalue}[1]{\gdef\@professor@value{#1}}
%</labreport|courseproj>
%    \end{macrocode}
% \end{macro}
%
% \begin{macro}{\professorinitials}\marg{initials}.
%
% Фамилия и инициалы преподавателя.
%
%    \begin{macrocode}
%<*labreport|courseproj>
\newcommand{\professorinitials}[1]{\gdef\@professor@initials{#1}}
%</labreport|courseproj>
%    \end{macrocode}
% \end{macro}
%
% \begin{macro}{\docyear}\marg{year}
%
% Год создания документа. По умолчанию \cs{year}.
%
%    \begin{macrocode}
%<*labreport|courseproj>
\gdef\@doc@year{\the\year}
\newcommand{\docyear}[1]{\gdef\@doc@year{#1}}
%</labreport|courseproj>
%    \end{macrocode}
% \end{macro}
%
% \subsection{Перечисления}
%
% \changes{v1.0e}{2025-02-14}{Устранение дублирования исходного текста во всех
%     файлах}
% \changes{v1.1}{2025-02-16}{Отказ от использования пакета enumitem}
% \textit{Из класса article.} Следующие команды используются, чтобы выставить 
% значения по умолчанию для параметров списка (среды list) (см. мануал \LaTeX\ 
% для объяснения значений этих параметров). Начальные установки таковы: 
% для начала правое поле списка \cs{rightmargin}, абзацный отступ 
% внутри списка \cs{listparindent} и отступ элемента \cs{itemindent} 
% устанавливаются в  0~пт. Затем для |K|-го уровня списка вызывается команда
% \cs{@listK}, где |K| обозначает римскую цифру от 1 до 6 (i, ii, ..., vi). 
% По умолчанию, \cs{@listK} должна устанавливать левое поле списка \cs{leftmargin} в 
% значение \cs{leftmarginK}.
%
% \ToDo{сделать настройку для дополнительных отступов}
%
% \begin{macro}{\leftmargini,\leftmarginii,\leftmarginiii,\leftmarginiv,\leftmarginv,\leftmarginvi,\itemindent,\listparindent}
% В ГОСТ предусмотрено, что элементы списков ведут себя как обычные абзацы:
% начинаются они с абзацного отступа, затем за ними идёт маркер или номер, а далее
% сам текст элемента. Никаких отступов от левого края страницы не предусмотрено.
%
%    \begin{macrocode}
%<*class>
\setlength\leftmargini{0pt}
\setlength\leftmarginii{0pt}
\setlength\leftmarginiii{0pt}
\setlength\leftmarginiv{0pt}
\setlength\leftmarginv{0pt}
\setlength\leftmarginvi{0pt}
%    \end{macrocode}
% \end{macro}
%
% \begin{macro}{\labelsep,\labelwidth}
% Между маркером/номером и текстом элемента списка стоит пробел, сам маркер/номер
% должен иметь естественную ширину. Пробел определим позднее.
%
%    \begin{macrocode}
\setlength\labelwidth{0pt}
\setlength\labelsep{0pt}
%    \end{macrocode}
% \end{macro}
%
% \begin{macro}{\@listi,\@listii,\@listiii,\@listiv,\@listv,\@listvi}
% Никаких дополнительных межстрочных отступов не предусмотрено.
%    \begin{macrocode}
\def\@listI{\leftmargin=\leftmargini
            \itemindent\parindent
            \listparindent\parindent
            \itemsep=0pt
            \parsep=0pt
            \topsep=0pt}
\let\@listi\@listI
\def\@listii{\leftmargin=\leftmarginii
            \itemindent\parindent
            \listparindent\parindent
            \itemsep=0pt
            \parsep=0pt
            \topsep=0pt}
\def\@listiii{\leftmargin=\leftmarginiii
            \itemindent\parindent
            \listparindent\parindent
            \itemsep=0pt
            \parsep=0pt
            \topsep=0pt}
\def\@listiv{\leftmargin=\leftmarginiv
            \itemindent\parindent
            \listparindent\parindent
            \itemsep=0pt
            \parsep=0pt
            \topsep=0pt}
\def\@listv{\leftmargin=\leftmarginv
            \itemindent\parindent
            \listparindent\parindent
            \itemsep=0pt
            \parsep=0pt
            \topsep=0pt}
\def\@listvi{\leftmargin=\leftmarginvi
            \itemindent\parindent
            \listparindent\parindent
            \itemsep=0pt
            \parsep=0pt
            \topsep=0pt}
%    \end{macrocode}
% \end{macro}
%
% \begin{macro}{\gostasbuk}\oarg{counter}
% Представляет значение счётчика в виде буквы русского алфавита, за
% исключением ё, з, й, о, ч, ъ, ы, ь.
%    \begin{macrocode}
\def\gostasbuk@alph#1{\ifcase#1\or
   а\or б\or в\or г\or д\or е\or ж\or
   и\or к\or л\or м\or н\or п\or р\or
   с\or т\or у\or ф\or х\or ц\or ш\or 
   щ\or э\or ю\or я%
   \else\ClassError{omgtudoc-asoiu}{Invalid value of counter}\fi
}
\def\gostasbuk#1{\expandafter\gostasbuk@alph\csname c@#1\endcsname}
%    \end{macrocode}
% \end{macro}
%
% Далее определим перечисления, пункты которых начинаются по стандарту.
% Для этого сначала переопределим среды |itemize| и |enumerate|:
%
% \begin{environment}{itemize}
% Ненумерованный список.
% \changes{v1.1b}{2025-02-21}{Исправлены абзацные отступы у
% ненумерованных списков}
%    \begin{macrocode}
\def\itemize{%
  \ifnum \@itemdepth >\thr@@\@toodeep\else
    \advance\@itemdepth\@ne
    \edef\@itemitem{labelitem\romannumeral\the\@itemdepth}%
    \expandafter
    \list
      \csname\@itemitem\endcsname
      {\def\makelabel##1{##1~}}%
  \fi}
%    \end{macrocode}
% \end{environment}
% \begin{environment}{enumerate}
% \changes{v1.1b}{2025-02-21}{Исправлены абзацные отступы у
% нумерованных списков}
% Нумерованный список.
%    \begin{macrocode}
\def\enumerate{%
  \ifnum \@enumdepth >\thr@@\@toodeep\else
    \advance\@enumdepth\@one
    \edef\@enumctr{enum\romannumeral\the\@enumdepth}%
    \expandafter
    \list
      \csname label\@enumctr\endcsname
      {\usecounter\@enumctr\def\makelabel##1{##1~}}%
  \fi}
%    \end{macrocode}
% \end{environment}
%
% Далее определим маркеры перечислений.
% \begin{macro}{\labelitemi,\labelitemii,\labelitemiii,\labelitemiv}
% Маркеры для ненумерованных списков. В соответствии со стандартом, это
% должно быть тире.
%    \begin{macrocode}
\renewcommand\labelitemi{---}
\renewcommand\labelitemii{---}
\renewcommand\labelitemiii{---}
\renewcommand\labelitemiv{---}
%    \end{macrocode}
% \end{macro}
% \begin{macro}{\theenumi,\theenumii,\theenumiii,\theenumiv}
% Представление счётчиков для нумерованныых списков. По стандарту
% это~--- буквы русского алфавита.
%    \begin{macrocode}
\renewcommand\theenumi{\gostasbuk{enumi}}
\renewcommand\theenumii{\gostasbuk{enumii}}
\renewcommand\theenumiii{\gostasbuk{enumii}}
\renewcommand\theenumiv{\gostasbuk{enumii}}
%    \end{macrocode}
% \end{macro}
% \begin{macro}{\labelenumi,\labelenumii,\labelenumiii,\labelenumiv}
% Метки для нумерованных списков. После буквы или цифры ставится круглая
% скобка.
%    \begin{macrocode}
\renewcommand\labelenumi{\theenumi)}
\renewcommand\labelenumii{\theenumii)}
\renewcommand\labelenumiii{\theenumiii)}
\renewcommand\labelenumiv{\theenumiv)}
%</class>
%    \end{macrocode}
% \end{macro}
%
% \Finale
\endinput
% vi: ts=2 sw=2 expandtab tw=72 formatoptions+=t linebreak
